\newpage
\thispagestyle{empty}
\addcontentsline{toc}{section}{Eidesstattliche Erklärung}
\section*{Eidesstattliche Erklärung}
\thispagestyle{empty}

Ich versichere hiermit an Eides statt, dass ich die vorliegende Projektarbeit mit dem oben genannten Titel selbstständig und ohne unzulässige fremde Hilfe erbracht habe. Ich habe keine anderen als die angegebenen Quellen und Hilfsmittel benutzt sowie wörtliche und sinngemäße Zitate kenntlich gemacht. Die Arbeit hat in gleicher oder ähnlicher Form noch keiner Prüfungsbehörde vorgelegen.

\vspace{1\baselineskip}

\subsection*{Erklärung zur Nutzung generativer KI-Werkzeuge}

In Übereinstimmung mit der Leitlinie der Fachhochschule Dortmund zum Einsatz Künstlicher Intelligenz in der Lehre vom 14.\,02.\,2024 mache ich die Nutzung generativer KI-Werkzeuge bei der Erstellung dieser Arbeit transparent. Die folgenden Werkzeuge wurden in den genannten Funktionen eingesetzt; alle übernommenen Inhalte wurden inhaltlich geprüft, eigenständig eingeordnet und gegen die zitierten Quellen abgesichert. Die wissenschaftliche Verantwortung für den Inhalt der Arbeit liegt vollständig bei mir.

\begin{table}[H]
\centering
\begin{tabular}{|p{4.0cm}|p{4.5cm}|p{4.5cm}|}
\hline
\textbf{Werkzeug (Version)} & \textbf{Einsatzzweck} & \textbf{Betroffene Arbeitsschritte} \\
\hline
% TODO (Igor): Modelle/Versionen pruefen und ergaenzen, z.B. ChatGPT (GPT-5), Claude (Modellversion)
ChatGPT / Claude (Anthropic) & Sprachliche Überarbeitung, Formulierungshilfen, Strukturvorschläge, kritische Durchsicht & Ausformulieren und Überarbeiten einzelner Absätze, Glättung von Übergängen, \LaTeX-Satz \\
\hline
% TODO (Igor): pruefen, ob so zutreffend
Claude (Anthropic) & Durchführung des byteweisen I/O-Vergleichs im dritten Durchlauf & Ausführung des Vergleichsskripts (Anhang~\ref{anhang:skript}) und Aufbereitung der Einzelergebnisse \\
\hline
% TODO (Igor): Werkzeug fuer Agentenprofile und Skills pruefen
Coding-Agenten (Claude, OpenCode) & Unterstützung bei der Erstellung der Pipeline-Konfiguration & Entwurf und Überarbeitung der Agentenprofile und Skills \\
\hline
OpenCode 1.14.20 bis 2.1.4 mit gpt-5.3-codex bzw.\ gpt-6.1-sol-ultrafast (OpenAI) & Ausführung der Migrations-Pipeline als Untersuchungsgegenstand & Erzeugung der Analyse-, Spezifikations- und Planungsartefakte sowie des Go-Zielcodes \\
\hline
NotebookLM (Google) & Recherche-Unterstützung und Abgleich von Aussagen mit den hinterlegten Quellen, anschließend manuell verifiziert & Vorrecherche und Quellenarbeit zu Themenkomplexen wie Spec-Driven Development und Reverse Engineering \\
\hline
\end{tabular}
\end{table}

Eingesetzte Prompts und Werkzeug-Konfigurationen wurden während der Bearbeitung dokumentiert und können auf Anforderung zur Einsichtnahme bereitgestellt werden. Nicht eingesetzt wurden generative KI-Werkzeuge zur Erzeugung wissenschaftlicher Aussagen oder zur Generierung von Zitaten und Literaturangaben ohne anschließende manuelle Verifikation.

\vspace{4\baselineskip}
Dortmund, \today \hfill \rule{6cm}{0.4pt}\\
\hspace*{0pt}\hfill Igor Besel
